\documentclass[11pt]{book}
\usepackage[T1]{fontenc}
\usepackage{lmodern}
\usepackage{amsmath,amsthm,amssymb}
\usepackage{fancyhdr}

\pagestyle{fancy}
\fancyhf{}
\fancyhead[LE,RO]{\thepage}
\fancyhead[RE]{\leftmark}
\fancyhead[LO]{\rightmark}

\theoremstyle{definition}
\newtheorem{definition}{Definition}[chapter]
\theoremstyle{plain}
\newtheorem{theorem}[definition]{Theorem}
\newtheorem{lemma}[definition]{Lemma}
\theoremstyle{remark}
\newtheorem*{remark}{Remark}

\title{A Gentle Introduction to Number Theory}
\author{Margaret Lindqvist}

\begin{document}
\frontmatter
\maketitle
\tableofcontents

\mainmatter
\chapter{Divisibility}

\section{The division algorithm}

\begin{definition}
Let $a,b\in\mathbb{Z}$. We say $a$ \emph{divides} $b$, written $a\mid b$, if there
exists $k\in\mathbb{Z}$ with $b = ka$.
\end{definition}

\begin{theorem}[Division algorithm]\label{thm:div}
For all $a\in\mathbb{Z}$ and $b\in\mathbb{Z}_{>0}$ there exist unique integers $q,r$
with $a = qb + r$ and $0\le r<b$.
\end{theorem}

\begin{proof}
Let $S = \{a - kb : k\in\mathbb{Z}\}\cap\mathbb{Z}_{\ge0}$. $S$ is non-empty, so by
the well-ordering principle it has a least element $r = a - qb$. If $r\ge b$ then
$r-b\in S$ is smaller, a contradiction. Uniqueness follows by subtracting two
representations.
\end{proof}

\section{Greatest common divisors}

\begin{lemma}\label{lem:gcd}
If $a = qb + r$ then $\gcd(a,b) = \gcd(b,r)$.
\end{lemma}

Repeated application of Lemma~\ref{lem:gcd} yields the Euclidean algorithm. For
example,
\begin{align*}
252 &= 1\cdot 198 + 54,\\
198 &= 3\cdot 54 + 36,\\
54  &= 1\cdot 36 + 18,\\
36  &= 2\cdot 18,
\end{align*}
so $\gcd(252,198) = 18$.

\begin{remark}
The number of steps is at most about five times the number of decimal digits of the
smaller argument (Lamé's theorem).
\end{remark}

\chapter{Primes}
\section{The fundamental theorem of arithmetic}
\begin{theorem}
Every integer $n\ge2$ is a product of primes, uniquely up to order.
\end{theorem}

\section{Infinitely many primes}
\begin{theorem}[Euclid]
There are infinitely many primes.
\end{theorem}
\begin{proof}
Suppose $p_1,\dots,p_n$ are all the primes and let $N = p_1p_2\cdots p_n + 1$. No $p_i$
divides $N$, contradicting Theorem~\ref{thm:div} applied to its prime factorisation.
\end{proof}

\backmatter
\chapter{Further Reading}
Hardy and Wright, \emph{An Introduction to the Theory of Numbers}, remains the classic
reference.

\end{document}
